\section{Compute Lifecycle：Pre-training、Post-training 与 Inference}

课堂反对“pre-training 已经结束”的二元叙事。Pre-training 高效吸收广泛数据并建立 base capability；post-training 把行为、工具和安全目标聚焦；inference compute 通过 context、search、reasoning 或多次采样解决具体任务。三者的边际收益、latency、cost 和可控性不同。

\subsection{不同阶段解决不同问题}

Pre-training 适合摊销到所有请求的广泛能力；post-training 适合产品/政策行为；inference 适合按任务动态投入。若 base model 缺少知识或 representation，增加 inference token 可能只是更长地猜；若 post-training 不足，强 base model 也可能不遵循工具 contract；若 serving budget 太低，复杂 reasoning 又无法落地。

\lecturefigure{15-compute-lifecycle.png}{Pre-training、post-training、inference compute 与 product feedback 形成互补生命周期。}{本地字幕 34:40--36:50；概念重绘。}

\begin{knowledgebox}{读图：比较 total cost，不只看训练 FLOPs}
一个模型的经济性包括训练摊销、post-training iteration、每请求 accelerator/time、context storage、cache、tool call 和 failure retry。更小 base model 加更多 inference search 可能适合低并发高价值任务；更强 base model 可能降低每请求步骤。选择应基于 target workload 与 SLO，而非阶段流行度。
\end{knowledgebox}

\subsection{本章小结}

Frontier progress 来自 compute 在生命周期中的重新分配，而不是某一阶段永久取代另一阶段。Program 要用端到端 quality、cost、latency 和 risk 选择组合。

